\begin{multicols*}{2}

\subsection*{Getting Started}

Find a Referee, create a knight, decide if they are a mage, a witch, or neither, and go for it.

\subsection*{Dice Rolls}

To resolve some events, a roll is needed.
Two mechanics modify the roll itself or the dice in the roll.

\paragraph{Upface \& Downface.}

Certain abilities can \texttt{Upface} and \texttt{Downface} a die.
To \texttt{Upface} a die, find it in the sequence:

\colbox{
    \[\texttt{d1, d4, d6, d8, d12}\]
}

And replace it with the next one. To \texttt{Downface} a die, replace it with the previous one.
\texttt{d1} can't be \texttt{Downfaced} and \texttt{d12} can't be \texttt{Upfaced}.

You probably don't have a \texttt{d1} in your possession, but as it always rolls \texttt{1}, you can use any small object to represent it.

When you are \texttt{Upfacing} or \texttt{Downfacing} a die in an \texttt{Attack} roll, do that before declaring \texttt{Focus} usage.

\newcolumn
\paragraph{Repeated advantage \& disadvantage.}

There are several types of rolls: \texttt{attack roll}, \texttt{Virtue Saves}, \texttt{Spellweaving Roll}, and more.
Some of them, including the named three, can be rolled with \texttt{Advantage} or \texttt{Disadvantage}.

To roll with \texttt{Advantage}, add a copy of a die that already exists in the roll.
Perform a roll and then remove a die that shares the number of faces with the added one.
Typically, you want to remove the lowest one, but ultimately it's for you to decide which one to remove.

It is possible to roll with repeated \texttt{Advantage} by adding and then removing the corresponding number of dice copies.

To roll with \texttt{Disadvantage} -- do the same, but you must remove the best one (the least one for \texttt{Virtue Saves}, the highest for \texttt{Attack Rolls}) from the roll.

If there is a need to simultaneously roll with \texttt{Advantage} and \texttt{Disadvantage}, remove matching pairs of "-antages" until only one of them remains.

It is possible to add \texttt{(Dis)Advantage} to rolls after the roll is performed.
To do that, choose a die that is present there and roll its copy.
Then proceed as if they were rolled together.


\end{multicols*}
